\documentclass[aspectratio=169,usenames,dvipsnames]{beamer}
\beamertemplatenavigationsymbolsempty
% \documentclass{beamer}

% UNCOMMENT FOR HANDOUTS
% \usepackage{pgfpages}
% \pgfpagesuselayout{2 on 1}[landscape,letterpaper,border shrink=5mm]
% \setbeameroption{show notes}

\mode<presentation>
\usetheme{Boadilla}
\usecolortheme{spruce}
\usepackage{../../LizStyle,ulem}
\usepackage{color}
\usepackage[color,matrix,arrow]{xy}
% \input xy
% \xyoption{all}
\usepackage{tikz}
\usepackage{tikz-cd}
\tikzset{commutative diagrams/diagrams={ampersand replacement=\&}}

\AtBeginSection{\frame{\sectionpage}}


% ==========Command for putting comments and removing them=======
% Visible:
\newcommand{\InstructorNotes}[1]{\textit{\color{orange}#1}}
\newcommand{\InstructorList}[1]{
\begin{itemize}
	\color{orange}
	#1
\end{itemize}
}
% Removed:
% \renewcommand{\InstructorNotes}[1]{}
% \renewcommand{\InstructorList}[1]{}


\newcommand{\bvfill}{\vskip0pt plus 1filll}

\usepackage{multimedia}
\usepackage{graphbox}

% \usepackage{algpseudocode}


\newcommand{\Reeb}{\textbf{Reeb}}
\newcommand{\Rtopc}{\textbf{Rtopc}}
% \newcommand{\RR}{\mathcal{R}}

\newcommand{\Set}{\textbf{Set}}
\newcommand{\Vect}{\textbf{Vect}}
\newcommand{\Top}{\textbf{Top}}
\newcommand{\Int}{\textbf{Int}}


\graphicspath{ {Figures/} {../Figures/}{../../Figures/} }



\begin{document}
%\pdfoptionpdfminorversion 5
%\logo{\includegraphics[height=1.5cm]{dukeshield.pdf}}
\title[]{Ch 7.2-7.4: Step Functions, Basis Functions, Start Splines }
\subtitle{Lecture 22 - CMSE 381}
\date{Mon, Oct 27, 2025}



\institute[MSU-CMSE]{Michigan State University\\
			::\\
          Dept of Computational Mathematics, Science \& Engineering}
\author[Dr.~Zhang]{Prof.~Mengsen Zhang}

\begin{frame}
 \titlepage


\end{frame}
%--------------------------
\begin{frame}[c]{Announcements}
\InstructorNotes{I split this class in two because it was too much last time, but now this is way too short. Or maybe I was just tired and skipping too much. Perhaps they should actually code the classification example? }
\begin{columns}
\column{.45\textwidth}
\centering

	\textbf{Last time:}
	\begin{itemize}
	\item 7.1 Polynomial regression 
	\item 7.2 Step functions
	\end{itemize}

	\vspace{.2in}

	\textbf{This lecture:}
	\begin{itemize}
	\item 7.2 Step functions
	\item 7.3 Basis functions
	\item 7.4 Regression Splines (Finish next lecture)
\end{itemize}


\column{.45\textwidth}
\centering

\textbf{Announcements:}
	\begin{itemize}
		\item HW \#6 Due Sunday
		\item Exam 2 graded. What worked? \href{https://PollEv.com/multiple_choice_polls/9vANQyoPcWiwin31KwCko/respond}{PollEv}
	\end{itemize}
\end{columns}

\end{frame}
%--- Next Frame ---%



%=================
\section{Last time}
%=================


\begin{frame}[c]{Polynomial regression}

\begin{columns}
\column{.45\textwidth}
\centering
Replace linear model
\begin{equation*}
	y_i = \beta_0 + \beta_1x_1 + \e_i
\end{equation*}
with
\begin{equation*}
	y_i = \beta_0 + \beta_1x_1 + \beta_2x_i^2 + \cdots + \beta_dx_i^d + \e_i
\end{equation*}

\column{.45\textwidth}
\centering
\InstructorList{
  \item Can learn with linear models by passing in predictors $x_i^\ell$
	\item Tend to not go higher than degree 3 or 4 because makes overly flexible
}

\end{columns}

\end{frame}
%--- Next Frame ---%

\begin{frame}[c]{Example with wage data}

\begin{columns}
\column{.4\textwidth}
\centering
\includegraphics[width = \textwidth, align = c]{Fig7-1a}

\column{.45\textwidth}
\centering
\InstructorList{
  \item Plot of wage vs age for men in central Atlantic region of the US
	\item Dark line is degree 4 polynomial
	\item Have variance for each coefficient
	\item Assume you have a model $\hat f(x_0) = \hat \beta_0 + \hat \beta_1 x_0 + \cdots + \hat \beta_d x_0^d$
	\item Can use that to estimate the pointwise variance $\Var(\hat f(x_0))$
	\item Draw 2 std deviations away from the line, 95\% confidence interval
}

\end{columns}
  $$
  -184.1542 + 21.24552 * age + -0.56386 * age^2 + 0.00681 * age^3 + -3e-05 * age^4
  $$

\end{frame}
%--- Next Frame ---%



\begin{frame}[t]{Step functions}
% \begin{equation*}
% 	\mathrm{I}(X < c) \qquad \mathrm{I}(c_1 \leq X < c_2) \qquad \mathrm{I} (c \leq X)
% \end{equation*}
\begin{columns}
\column{.25\textwidth}
\centering
\begin{itemize}
	\item $\mathrm{I}(X<c)$
	\item $\mathrm{I}(c_1 \leq X <c_2)$
	\item $\mathrm{I}(c \leq X)$
\end{itemize}


\column{.45\textwidth}
\centering
\includegraphics[width = \textwidth, align = c]{Eq7-4}
\end{columns}
	
\bvfill 
\centering 
Learned model: 
\begin{equation*}
	y_i = \beta_0 + \beta_1C_1(x_i) + \beta_2 C_2(x_i) + \cdots + \beta_K C_K(x_i) + \e_i
\end{equation*}

\bvfill 
\end{frame}
%--- Next Frame ---%

\begin{frame}[t]{Example: Cut points at -4, -1, 3, 6}
\begin{columns}
\column{.3\textwidth}
\centering
$C_0(X)$
\includegraphics[align = c, width = \textwidth]{StepFunctionExample-C0}





\column{.3\textwidth}
\centering
$C_1(X)$
\includegraphics[align = c, width = \textwidth]{StepFunctionExample-C1}

\vspace{.2in}

$C_2(X)$
\includegraphics[align = c, width = \textwidth]{StepFunctionExample-C2}

\vspace{.2in}

$C_3(X)$
\includegraphics[align = c, width = \textwidth]{StepFunctionExample-C3}

\column{.3\textwidth}
\centering
$C_4(X)$
\includegraphics[align = c, width = \textwidth]{StepFunctionExample-C4}

\end{columns}
\end{frame}
%--- Next Frame ---%
% {
% \setbeamercolor{background canvas}{bg=Apricot!25}
% \begin{frame}[t]{Coding with step functions}
% \begin{columns}
% \column{.45\textwidth}
% \centering
% \column{.45\textwidth}
% \centering
% \end{columns}
% \end{frame}
% }
% %--- Next Frame ---%


\begin{frame}[c]{Step function example}

\begin{columns}
\column{.45\textwidth}
\centering
\includegraphics[width = \textwidth, align = c]{Fig7-2a}

\column{.45\textwidth}
\centering
\InstructorList{
  \item Learned peak where the high earners are

}

\end{columns}

\end{frame}
%--- Next Frame ---%
\begin{frame}[t]{What will you learn today?}
\begin{itemize}
	\item How to fit step function models for classification problems?
		\begin{itemize}
			\item You should also be able to implement them in \texttt{Python}.
		\end{itemize}
	\item What are basis functions?
		\begin{itemize}
			\item You should be able to articulate what the basis functions $b_i$ are for different models that we have covered.
			\item Examples include polynomial, step functions, and cubic splines.
		\end{itemize}
	\item What is the purpose of using basis functions?
		\begin{itemize}
			\item How do they influence model flexibility?
			\item How does this purpose manifest in different models, such as regression splines?
		\end{itemize}
	\item How to define cubic splines mathematically?
		\begin{itemize}
			\item You should be able to write down the equations for the model between knots.
			\item And the equations for constraints at the knots.
		\end{itemize}
	\item How to calculate the cubic spline coefficients by hand by applying the above mathematical definitions?
\end{itemize}
\end{frame}

%=================
\section{Classification versions}
%=================

\begin{frame}[t]{Remember logisitic regression?}
\begin{columns}
\column{.45\textwidth}
\centering

\includegraphics[width = \textwidth, align = c]{LogisticFunction}
\column{.45\textwidth}
\centering

\begin{equation*}
	p(X) =
	\frac{e^{\beta_0 + \beta_1X}}{1 +e^{\beta_0 + \beta_1X} }
\end{equation*}

\vspace{.5in}

\textbf{Multiple features:}
\begin{equation*}
	p(X) = \frac{e^{\beta_0 + \beta_1X_1 + \cdots + \beta_pX_p}}{1+e^{\beta_0 + \beta_1X_1 + \cdots + \beta_pX_p}}
\end{equation*}
\end{columns}
\end{frame}
%--- Next Frame ---%
\begin{frame}[c]{Classification version: Polynomial regression}

	\begin{equation*}
		\Pr(y_i>250 \mid x_i) =
		\frac{\exp(\beta_0 + \beta_1x_i + \cdots + \beta_dx_i^d)}
			 {1+\exp(\beta_0 + \beta_1x_i + \cdots + \beta_dx_i^d)}
	\end{equation*}
	\begin{columns}
	\column{.35\textwidth}
	\centering
	\includegraphics[width = \textwidth, align = c]{Fig7-1b}
	\column{.45\textwidth}
	\centering
	
	\footnotesize
	\InstructorList{
	  \item Note in previous fig that there is a distinct subcluster of high earners making more than \$250K
		\item Build a logistic regression model as above
		\item Note that the 95\% confidence interval gets very wide on the right side
		\item Large sample size $n=3,000$ but small number (79) of high earners makes for high variance in coeffs and wide confidence intervals
	
	}
	
	\end{columns}
	
	\end{frame}
	%--- Next Frame ---%
\begin{frame}[c]{Classification version: Step functions }
\begin{equation*}
	\Pr(y_i>250 \mid x_i) =
	\frac{\exp(\beta_0 + \beta_1C_1(x_i) + \beta_2 C_2(x_i) + \cdots + \beta_K C_K(x_i))}
	     {1+\exp(\beta_0 + \beta_1C_1(x_i) + \beta_2 C_2(x_i) + \cdots + \beta_K C_K(x_i))}
\end{equation*}
\begin{columns}
\column{.35\textwidth}
\centering
\includegraphics[width = \textwidth, align = c]{Fig7-2b}
\column{.45\textwidth}
\centering
\InstructorList{
  \item Again, learns that middle bit
	\item Still big confidence interval on the right, but likely because there's just less data over there
}

\end{columns}

\end{frame}
%--- Next Frame ---%

{
\setbeamercolor{background canvas}{bg=Apricot!25}
\begin{frame}[t]{Coding bit: classification version}
\begin{columns}
\column{.45\textwidth}
\centering
\includegraphics[align = c, width = \textwidth]{WageColoredBy250.png}
\column{.45\textwidth}
\centering
\InstructorNotes{Just talk through on projector, there's nothing in there for them to code}
\end{columns}
\end{frame}
}
%--- Next Frame ---%

% \begin{frame}[t]{Computation in the jupyter notebook}
% \begin{columns}
% \column{.35\textwidth}
% \centering
% \includegraphics[align = c, width = \textwidth]{WageColoredBy250.png}

% \column{.55\textwidth}
% \centering
% \includegraphics[align = c, width = \textwidth]{StepFunctions-codescreenshot}

% \end{columns}
% \end{frame}
%--- Next Frame ---%
\begin{frame}[c]{A few more comments on step functions}

\begin{columns}
\column{.45\textwidth}
\centering
\begin{itemize}
  \item Gives the chance to break up the domain, avoid forcing global structure
	\item Need to make decisions about the $c_i$.  A bit arbitrary unless your data has natural breakpoints.
	\item Popular in biostats and epidemiology
\end{itemize}

\column{.45\textwidth}
\centering

\end{columns}

\end{frame}
%--- Next Frame ---%

%=================
\section{Basis functions}
%=================


\begin{frame}[t]{Basis Functions Setup}
Polynomial and piecewise-constant regression models are special cases of a \textit{basis function} approach.

\begin{equation*}
	y_i = \beta_0 + \beta_1b_1(x_i) + \beta_2 b_2(x_i) + \cdots + \beta_K b_K(x_i) + \e_i
\end{equation*}

\begin{columns}
\column{.45\textwidth}
\centering
\InstructorList{
  \item Pick a collection of basis functions $b_1(X),\cdots,b_K(X)$
	\item Use least squares to figure out the coefficients
	\item Explain the $b_i$'s for polynomial and stepwise functions
	\item incontrast to the dim-reduction methods, here to want to add dimensions. 
	% \item Get to still use the inference tools from Ch 3 like standard errors and F-stats

}

\column{.45\textwidth}
\centering

\InstructorList{
	\item Lots of possible options for these
	\InstructorList{
	  \item Some examples are wavelets or Fourier series
		\item Next section is we'll talk about regression splines
	}
}


\end{columns}

\end{frame}
%--- Next Frame ---%


%==========
\section{Regression Splines}
%==========


\begin{frame}[c]{Piecewise polynomials}

\begin{columns}
\column{.45\textwidth}
\centering
\begin{itemize}
	\item Fit a polynomial regression
\begin{equation*}
	y_i = \beta_0 + \beta_1x_1 + \beta_2x_i^2 + \cdots + \beta_dx_i^d + \e_i
\end{equation*}
\item Let the $\beta_i$'s be different at different locations of the range.
\end{itemize}
\column{.45\textwidth}
\centering
\InstructorList{
  \item Points where the coefficients change are called "knots"
	\item A piecewise cubic with no knots is just standard cubic polynomial
	\item More knots gives more flexible polynomial
	\item Piecewise constant of previous section is degree 0 version of this
}

\end{columns}

\end{frame}
%--- Next Frame ---%

\begin{frame}[c]{Example of piecewise polynomial}

\begin{columns}
\column{.45\textwidth}
\centering
\includegraphics[width = \textwidth, align = c]{Fig7-3a}
\column{.45\textwidth}
\centering

Example:
\includegraphics[width = \textwidth, align = c]{Eq7-8b.png}
\end{columns}

\InstructorList{
\item Uses 8 degrees of freedom to fit the model
  \item problem: we have a discontinuous function.

}


\end{frame}
%--- Next Frame ---%

\begin{frame}[c]{The fix}

\begin{columns}
\column{.45\textwidth}
\centering
\begin{itemize}
	\item Fit piecewise polynomial
	\item Require continuity at knots
\end{itemize}
\column{.45\textwidth}
\centering
\includegraphics[width = \textwidth, align = c]{Fig7-3b}
\end{columns}
\InstructorNotes{Looks better, but pointy join is odd}
\end{frame}
%--- Next Frame ---%

\begin{frame}[t]{The better fix: Cubic splines}

\begin{columns}
\column{.45\textwidth}
\centering

\begin{itemize}
	\item Fit piecewise polynomial
	\item Require continuity at knots
	\item Require the first and second derivatives to be continuous at knots
\end{itemize}
\InstructorList{
  \item Now we have smooth joins
	\item This is called a "cubic spline"
	\item 8 original degrees of freedom, plus three constraints (continuous, 1st deriv continuous, 2nd deriv continuous) leaves 5 degrees of freedom
	\item cubic spline with $K$ knots uses $4+K$ degrees of freedom
}
\column{.45\textwidth}
\centering
\includegraphics[width = \textwidth, align = c]{Fig7-3c}

Test your understanding: \href{https://PollEv.com/multiple_choice_polls/bkj0zshpQHzhiwM3UXroS/respond}{PollEv}
\end{columns}

\end{frame}
%--- Next Frame ---%
{
\setbeamercolor{background canvas}{bg=Apricot!25}
\begin{frame}[t]{Example}
\begin{columns}
\column{.45\textwidth}
\centering
We have the following piecewise cubic polynomial: 
\begin{equation*}
	f(x) = 
	\begin{cases}
		2+x+x^{2}+0.1x^{3} & x \leq 1 \\
		b_0 + b_1 x + b_2 x^2 -x^3 & x > 1
	\end{cases}
\end{equation*}
What are $b_1$, $b_1$, and $b_2$ to make this a cubic spline?
\column{.45\textwidth}
\centering
\includegraphics[align = c, width = \textwidth]{SplineMatchingExample.png}

\end{columns}

\bvfill
 Check your answers: \href{https://www.desmos.com/calculator/kbm0zivqco}
{www.desmos.com/calculator/kbm0zivqco}
\end{frame}
}
%--- Next Frame ---%
{
\setbeamercolor{background canvas}{bg=Apricot!25}
\begin{frame}[t]{More space for work}

\begin{columns}
\column{.45\textwidth}
\centering
	\begin{equation*}
		f(x) = 
		\begin{cases}
			2+x+x^{2}+0.1x^{3} & x \leq 1 \\
			b_0 + b_1 x + b_2 x^2 -x^3 & x > 1
		\end{cases}
	\end{equation*}
	\InstructorNotes{
Instructor tool for generating solution while messing with $b_3$: \url{https://www.desmos.com/calculator/5prmps9cuo}
	}
\column{.45\textwidth}

\centering
\InstructorNotes{
	\begin{equation*}
		f(x) = 
		\begin{cases}
			2+x+x^{2}+0.1x^{3} & x \leq 1 \\
			3.1 + -2.3 x + 4.3 x^2 -x^3 & x > 1
		\end{cases}
	\end{equation*}
\includegraphics[align = c, width = .6\textwidth]{SplineMatchingExample-Soln}
	}

\end{columns}
\end{frame}
}


\begin{frame}[t]{Next time}

\includegraphics[align = c, width = .45\textwidth]{Schedule-F25-3.png}

\InstructorNotes{Decided to skip 7.5-7.7}

\end{frame}
%--- Next Frame ---%

\end{document}
